\subsection{Finite geometric cancellation and every nonunit singleton at three outer edges}\label{proof:R2-C045}
\paragraph{R45.1 Original model and exact statements}

Work on the original finite probability space after deleting zero-mass atoms,
with a real skew kernel \(K\) satisfying \(|K(x,\allowbreak y)|\le1\).
All norms, inner products, adjoints and expectations use the same original measure.
Write
\fitdisplay{
Jv=(Ev)1,\quad A=J+K,\quad T=J-K=A^*,\quad
 g=K1,\quad p=1+g,\quad q=1-g,
}{}
\fitdisplay{
f_n=A^{n}\one,\quad \bar f_n=T^{n}\one,\quad
 h_n=T f_{n-1},\quad W_n=A \bar f_{n-1},\quad
 \delta_j=\|K^{j}\one\|_2^2,\quad d=\delta_1.
}{}
In particular p(x), \(q(x)\ge \mu(x)>0.\) A finite actual tree message is positive
because the all-equal labeling contributes positive original mass.
Define
\fitdisplay{
I_q(U)=E[q^{-1/3}|U|^{8/3}],\qquad B_n=E[\bar f_n^4/q],
}{}
and let \(M_v\) be multiplication by the actual function v. The supplied
finite calculus gives
\fitdisplay{
\theta=13/32,\quad\kappa=19/32,\quad
 0\le d\le1/3,\quad d^2\le \delta_2\le\theta d,\quad
 \delta_{j+1}\le\theta \delta_j,\quad \|K\|^2\le\theta.
}{R45.1.1}

Theorem A: exact improved bare-row bounds. For every integer \(n\ge 4,\allowbreak \)
\fitdisplay{
B_n\le
 \begin{cases}
 1-7d/4,&0\le d\le2/15,\\
 1-3d/2,&0\le d\le7/40,\\
 1-29d/20,&0\le d\le1/5.
 \end{cases}
}{R45.1.2}
The overlapping bounds are all valid on their stated domains for the
same actual n-edge path. Section3 also proves a single explicit continuous
bound on the entire original domain, before these convenient relaxations.
For \(n\ge 5\) there is the separate all-domain estimate
\fitdisplay{
B_n\le8(1-2d)(1-5d/2)^3.
}{R45.1.3}
The restriction \(n\ge 5\) in (R45.1.3) is essential to the present proof.

Theorem B: every nonunit singleton, outer length three. Suppose the
actual internal vertex has, after maximal genuine equal-length pairing,
exactly one residual arm of length \(a\ge 2.\) Its actual terminal has any
strict residual list of size at least two. Let \(L\ge 1\) be the inner original
connector length. Include every original root and terminal equal pair.
Section5 prescribes a direction depending only on these integers. For
outer length \(M=3,\allowbreak \) its actual message into the original center satisfies
\fitdisplay{
\boxed{I_q(U)^3\le e^{-d/80}.}
}{R45.1.4}
There is no relation required between a, L and the terminal lengths and
no ordinary-eighth input assumption. 

Theorem C: precise union and original P5 consequence. The
E008 construction and Theorem B together give, for every such half-tree,
\fitdisplay{
I_q(U)^3\le e^{-\eta_M d},\qquad
 \eta_M=
 \begin{cases}
 1/2,&M=1,\\
 1/80,&M=3,\\
 1/2+\lfloor(M-1)/3\rfloor/2,&M\ge4.
 \end{cases}
}{R45.1.5}
 A P5 with both noncentral internal residuals
nonunit singletons and both outer lengths different from two therefore
has an original structure-only orientation with \(Z\le 1\) for arbitrary
terminal and central multisets, and both independent positive inner
lengths. When both ends are nonreducible and the center has a nonempty
R residual, the common strict bound is \(Z\le \exp(-d/320).\) Sections~\ref{sec:five-chain} and~\ref{sec:completion} give
the comparison branches and the separately qualified mixed use.

\paragraph{R45.2 The finite geometric identity which sharpens the C100 row}

Define
\fitdisplay{
G_m(z)=\sum_{j=0}^{m-1}z^j,\qquad
 E_r(x)=\sum_{j=0}^{r-1}(-x)^j\quad(r\ge1).
}{}
For x in [0, theta], \(1-x\le E_r(x)\le 1.\) All sums here are finite.
At the original spectral block \(z=i \sqrt{x}\), write
\fitdisplay{
G_m(z)=A_m(x)+zD_m(x).
}{}
The letter \(D_m\) is a scalar polynomial here, not an orientation. Fix
\(n\ge 4\) and \(m=n-1\ge 3.\) In the actual C100 notation set
\fitdisplay{
R=K G_m(-K)g,\quad h=E[gR],\quad V=ER^2.
}{}
Let \(\nu_g\) be the actual finite spectral measure of -\(K^2\) at g. Its mass
is d and its first moment is \(\delta_2.\) The exact identities are
\fitdisplay{
h=\int xD_m(x)\,d\nu_g(x),\quad
 V=\int x\{A_m(x)^2+xD_m(x)^2\}\,d\nu_g(x).
}{R45.2.1}
This is the same actual R and h as in C100; it is not an independent
replacement vector. We now retain cancellation in V-h.

If \(m=2r\) is even, then \(r\ge 2\) and \(A_m=D_m=E_r.\) Hence
\fitdisplay{
x(A_m^2+xD_m^2-D_m)
 =(-1)^{r+1}x^{r+1}E_r.
}{R45.2.2}
Indeed (\(1+x)E_r=1-(-x)^r.\) When the right side is positive, r is odd
and \(r\ge 3.\) It is then at most \(x^4,\allowbreak \) which is at most
\(x^3(1-x+x^2),\allowbreak \) since \(1-x+x^2-x=(1-x)^2\ge 0.\) Negative cases already
satisfy the same upper bound.

If \(m=2r+1\) is odd, \(r\ge 1,\allowbreak \) then \(D_m=E_r\) and \(A_m=E_{r+1}=1-xE_r.\)
Direct substitution gives
\fitdisplay{
A_m^2+xD_m^2-D_m=(-x)^r A_m.
}{R45.2.3}
Only even \(r\ge 2\) gives a positive right side. In that case
\fitdisplay{
A_m=\frac{1+x^{r+1}}{1+x}
 \le\frac{1+x^3}{1+x}=1-x+x^2,
}{}
and \(x^{r+1}\le x^3.\) Thus every \(m\ge 3,\allowbreak \) including m3 and m4, obeys
\fitdisplay{
x(A_m^2+xD_m^2-D_m)\le x^3(1-x+x^2).
}{R45.2.4}
The derivative of \(x^2(1-x+x^2)\) is \(x(2-3x+4x^2)\ge 0:\) the quadratic
has negative discriminant and positive leading coefficient. Consequently
\fitdisplay{
\boxed{V-h\le r_*\delta_2,\qquad
 r_*=\theta^2(1-\theta+\theta^2)=\frac{131313}{1048576}.}
}{R45.2.5}
This replaces the weaker C100 bound \(V-h\le \theta \delta_2.\) The estimate is
uniform over all finite \(m\ge 3.\) It makes no assumption \(x\ge d\) and uses no
infinite geometric series. Useful exact scalar comparisons are
\fitdisplay{
\frac18<r_*<\frac{13}{100},\qquad
 \theta r_* =\frac{1707069}{33554432}<\frac{51}{1000}.
}{R45.2.6}
For the last comparison, \(1707069000<1711276032.\) Put
\fitdisplay{
\delta_*:=\theta^2-51/1000=\frac{14597}{128000}>\frac9{80}.
}{R45.2.7}
It follows that \(\theta(\theta-r_*)>\delta_*.\)

\paragraph{R45.3 Same-row fourth expansion and a continuous explicit bound}

This section reuses the exact C100 calculation only through its
given identities. To avoid confusing an arm length with a scalar,
put \(\alpha=\kappa\) d, \(C=1-\alpha,\allowbreak  t=1+d/2.\) The two positive finite-convolution
upper bounds on the SAME actual path are
\fitdisplay{
\bar f_n\le Cq+R-\alpha,\qquad \bar f_n\le2Cq.
}{R45.3.1}
They hold for every \(n\ge 4.\) The second retains the actual mean of \(\bar f_{n-2}.\)
For \(s\ge -1,\allowbreak \)
\fitdisplay{
\min(1+s,2)^4\le1+4s+7s^2+4s^3.
}{}
Applying this at \(s=(R-\alpha)/(Cq),\allowbreak \) then retaining the correlated cubic
through \(W=2R-Cg,\allowbreak \) gives the exact C100(2.10):
\fitdisplay{
\begin{aligned}
 B_n\le{}&C^4(1+3d)+C^3[td-6\alpha d-8\alpha]
                 +C\alpha^2[19-d-33\alpha]\\
 &+C B(d)(V-h)
  +C[-1+(13+d)\alpha-30\alpha^2]h-C^2E(gR^2),
\end{aligned}
}{R45.3.2}
where
\fitdisplay{
B(d)=9-\frac{481}{32}d>0.
}{}
Here B(d) is a scalar coefficient; it is unrelated to the index in \(B_n.\)
The beta monotonicity, cubic signs and coefficient collection behind
(R45.3.2) are unchanged Q5-C100 steps, given.

The genuine h bounds are kappa \(\delta_2\le h\le \delta_2.\) Separating the three terms
of its coefficient by sign, as C100 does, gives
\fitdisplay{
[-1+(13+d)\alpha-30\alpha^2]h
 \le H(d)\delta_2,
}{}
\fitdisplay{
H(d)=-\kappa+13\kappa d+\kappa(1-30\kappa^2)d^2.
}{R45.3.3}
After (R45.2.5), the combined coefficient of C \(\delta_2\) is now
\(S_new(d)=r_*B(d)+H(d).\) It is positive throughout [0, \(1/3\)]. Indeed
\fitdisplay{
9r_*-\kappa>17/32,\quad
 13\kappa-(481/32)r_*>5,\quad
 \kappa(1-30\kappa^2)=-93157/16384>-6,
}{}
so \(S_new(d)>17/32+5d-6d^2>0.\) Only AFTER this sign check may we
replace the combined \(\delta_2\) by theta d.

The exact old C100 polynomial was
\fitdisplay{
\begin{aligned}
 P(d)=1&-\frac{963}{512}d+\frac{171113}{32768}d^2
 -\frac{7275575}{1048576}d^3\\
 &+\frac{14103187}{16777216}d^4
 +\frac{1063145}{1048576}d^5.
\end{aligned}
}{}
The only changed term is therefore exactly
-C B(d) \(\theta(\theta-r_*)d.\) The unchanged actual correlation satisfies
\(|E(gR^2)|\le (2/3)d^{3/2}.\) Thus
\fitdisplay{
B_n\le P(d)+\frac23 C^2d^{3/2}
             -C B(d)\theta(\theta-r_*)d.
}{R45.3.4}
The already proved C100 rational comparison, valid on all \(d\le 1/3,\allowbreak \) is
\fitdisplay{
(P(d)-1)/d\le-963/512+(21/4)d-(32/5)d^2.
}{R45.3.5}
It follows from \(A2<21/4,\allowbreak  A3<-69/10,\allowbreak  A4<1,\allowbreak  A5<33/32\) and
\(d^3+(33/32)d^4\le (43/96)d^2\).

Combining (R45.2.7),(R45.3.4),(R45.3.5), for \(d>0\) we obtain the claimed continuous
same-row bound
\fitdisplay{
\boxed{B_n\le1-s(d)d,}
}{}
\fitdisplay{
s(d)=\frac{963}{512}-\frac{21}{4}d+\frac{32}{5}d^2
       -\frac23(1-\kappa d)^2\sqrt d
       +\delta_*(1-\kappa d)\left(9-\frac{481}{32}d\right).
}{R45.3.6}
It extends continuously to \(d=0.\) This is an analytic formula, not a
finite scan. It holds for every \(n\ge 4\) without changing the actual path.

The function s is decreasing on [0, \(1/3\)]. The first polynomial's derivative
is at most -\(59/60.\) The product C B(d) is decreasing with both factors
positive. Finally \(C^2 \sqrt{d}\) is increasing: for \(d>0\) its derivative has
the sign of C(1-5kappa d), and \(1-5\kappa/3=1/96>0.\) These three signs
prove monotonicity on the whole interval, including the continuous
endpoint at zero.

\paragraph{R45.4 Exact endpoint payments and the higher-domain row}

It now suffices to lower-bound s at three exact rational endpoints.
Evaluate the monotone analytic formula at those endpoints. Write \(F(d)=963/512-(21/4)d+(32/5)d^2.\)

\paragraph{R45.4.1 Endpoint \(d=2/15\)}

Here \(C=221/240>23/25\) and \(B=1679/240>69/10.\) Therefore
\fitdisplay{
\delta_*CB>\frac9{80}\frac{23}{25}\frac{69}{10}
 =\frac{14283}{20000}>7/10.
}{}
Also \(F(2/15)=963/512-7/10+128/1125>129/100:\) use
\(963/512>47/25\) and \(128/1125>11/100.\) Since
\(C^2=48841/57600<17/20\) and \(\sqrt{2/15}<11/30,\allowbreak \)
\fitdisplay{
\frac23 C^2\sqrt{2/15}<\frac{187}{900}.
}{}
Consequently
\fitdisplay{
s(2/15)>129/100+7/10-187/900=401/225>7/4.
}{R45.4.1}

\paragraph{R45.4.2 Endpoint \(d=7/40\)}

Here \(C=1147/1280>89/100\) and \(B=8153/1280>63/10.\) Hence
\fitdisplay{
\delta_*CB>\frac9{80}\frac{89}{100}\frac{63}{10}
 =\frac{50463}{80000}>5/8.
}{}
The exact F is \(963/512-147/160+49/250>23/20,\allowbreak \) by
\(963/512>47/25,\allowbreak  147/160<23/25,\allowbreak  49/250>19/100.\)
Furthermore \(C<9/10\) and \(\sqrt{7/40}<21/50,\allowbreak \) so
\fitdisplay{
\frac23 C^2\sqrt{7/40}<\frac{567}{2500}<11/40.
}{}
It follows that
\fitdisplay{
s(7/40)>23/20+5/8-567/2500>3/2.
}{R45.4.2}

\paragraph{R45.4.3 Endpoint \(d=1/5\)}

Now \(C=141/160,\allowbreak  B=959/160\) and \(CB=135219/25600.\) The retained gain
is greater than \(3/5,\allowbreak \) because
\fitdisplay{
14597\cdot135219=1973791743>1966080000
 =\frac35(128000\cdot25600).
}{}
The remaining exact comparisons are
\fitdisplay{
F(1/5)=\frac{69559}{64000},\quad
 \sqrt{1/5}<9/20,\quad
 \frac23C^2\sqrt{1/5}<\frac{59643}{256000}.
}{}
Thus
\fitdisplay{
s(1/5)>\frac{372193}{256000}
          >\frac{371200}{256000}=29/20.
}{R45.4.3}
Equations (R45.4.1)-(R45.4.3) and monotonicity prove every line of (R45.1.2).

\paragraph{R45.4.4 The original finite path-mean bound}

For clarity, the exact finite argument used on the other domain
is included. Define \(\beta_1=\beta_2=0\) and
\fitdisplay{
\beta_j=\sum_{r=1}^{\lfloor(j-1)/2\rfloor}(-1)^{r-1}\delta_r,
 \qquad \tau=d-\delta_2.
}{}
The actual finite positive convolution is
\fitdisplay{
\sum_{r=0}^n(-K)^r1=\bar f_n+\sum_{j=3}^n\beta_j\bar f_{n-j},
 \qquad \tau\le\beta_j\le d\quad(j\ge3).
}{R45.4.4}
Its exact means \(\mu_n=E\bar f_n=Ef_n\) satisfy
\fitdisplay{
\begin{aligned}
 \mu_2&=1-d,&\mu_3&=1-2d,\\
 \mu_4&=1-3d+\delta_2,&\mu_5&=1-4d+2\delta_2+d^2,\\
 \mu_6&=1-5d+3\delta_2+3d^2-\delta_3,\\
 \mu_7&=1-6d+4\delta_2+6d^2-2\delta_3-2d\delta_2.
\end{aligned}
}{}
For n4, 5, 6, 7 their losses divided by d are at least respectively
\(83/32,\allowbreak  137/48,\allowbreak  89/32,\allowbreak  127/48.\) These follow from \(d\le 1/3\) and
\(\delta_2\le \theta\) d, retaining the displayed positive losses. For example,
\fitdisplay{
6d-(4-2d)\delta_2-6d^2
 \ge[6-4\theta-(6-2\theta)d]d\ge127d/48.
}{}
For \(n\ge 8,\allowbreak \) take the expectation in (R45.4.4) and retain the six actual
terms corresponding to \(\mu_0,\allowbreak ...,\allowbreak  \mu_5,\allowbreak \) together with \(\beta_{n+1}.\)
The three coefficients \(\beta_{n+1},\allowbreak  \beta_n,\allowbreak  \beta_{n-1}\) are at least
tau+kappa \(\delta_3.\) This gives
\fitdisplay{
\begin{aligned}
 1-\mu_n
 &\ge\tau(7-10d+3\delta_2+d^2)+3\kappa \delta_3\\
 &\ge\kappa d(7-10d+7d^2)\ge95d/36>83d/32.
\end{aligned}
}{R45.4.5}
For the middle step use \(\delta_2\ge d^2\) and \(\delta_3\ge d \delta_2\ge d^3\); the latter
is the Cauchy moment inequality \(\delta_2^2\le d \delta_3\) combined with \(\delta_2\ge d^2.\)
The bracket is positive. The final quadratic decreases on [0, \(1/3\)].
Hence
\fitdisplay{
\mu_j\le1-2d\ (j\ge3),\qquad
 \mu_j\le1-83d/32\le1-5d/2\ (j\ge4).
}{R45.4.6}


If \(n\ge 5,\allowbreak \) the positive original kernel gives \(Tv\le 2Ev\) and
\(T^2v\le 2(Ev)q\) for \(v\ge 0.\) Thus, for the SAME path,
\fitdisplay{
\bar f_n\le2(1-2d)q,
 \quad \bar f_n\le2(1-5d/2),\quad E\bar f_n\le1-5d/2.
}{}
All right sides are nonnegative for \(d\le 1/3.\) Consequently
\fitdisplay{
B_n\le2(1-2d)E \bar f_n^3
 \le8(1-2d)(1-5d/2)^3,
}{}
which is (R45.1.3). The first mean is \(\mu_{n-2}\); it is exactly mu3 when
\(n=5.\) The latter two require indices at least four. No n4 endpoint is
hidden in this argument.

\paragraph{R45.5 Every actual terminal and the original weighted input}

Theorem B uses the E008 Section10 construction. The
complete outward-reference table is reproduced to make its integer-only
choice explicit. Canonical pairs mean \(h_s f_t\) with \(s<t\); \(h_{1}=q.\) Every
genuine terminal equal-pair decoration stays inside its original
connector. The full outward connector is \(T^L.\)


\begingroup\small
\setlength{\tabcolsep}{3pt}
\setlength{\LTleft}{0pt}\setlength{\LTright}{\fill}
\begin{longtable}{@{}p{0.2200\linewidth}p{0.7000\linewidth}@{}}
\caption{Finite geometric cancellation and every nonunit singleton at three outer edges: parameter values.}\label{tab:parameters:r2-c045:1}\\
\toprule
Entry & Actual outward-reference \(F_ref\) \\
\midrule\endfirsthead
\toprule
Entry & Actual outward-reference \(F_ref\) \\
\midrule\endhead
\bottomrule\endfoot
Exact (1, b) & \(q f_b\). \\
\addlinespace[2pt]
Exact (2, b) & \(h_{2} f_b\). \\
\addlinespace[2pt]
Exact (3, b), \(L=1\) or2 & \(f_{3} \bar f_b\). \\
\addlinespace[2pt]
Exact (3, b), \(L\ge 3\) & \(h_{3} f_b\). \\
\addlinespace[2pt]
Exact pair of \(minimum\ge 4\); every even multipair & Product of all canonical \(h_s f_t\). \\
\addlinespace[2pt]
Odd \(minimum\ge 2,\allowbreak \) second \(length\ge 4\) & \(\bar f_a\) times all later canonical pairs. \\
\addlinespace[2pt]
Odd prefix (2, 3, c) & \(\bar f_2 f_{3} h_c times later canonical pairs\). \\
\addlinespace[2pt]
Exact (1, s, t), \(s\ge 4\) & \(q f_s h_t\). \\
\addlinespace[2pt]
Unit-minimum odd, \(second\ge 4,\allowbreak \) at least five arms & q times all later canonical pairs. \\
\addlinespace[2pt]
Unit-minimum odd prefix (1, 3, c) & \(q f_{3} h_c times later canonical pairs\). \\
\addlinespace[2pt]
Exact (1, 2, 3), \(L=1\) & \(qf_2\bar f_3\). \\
\addlinespace[2pt]
Exact (1, 2, 3), \(L\ge 2\) & \(qf_2h_3\). \\
\addlinespace[2pt]
Exact (1, 2, c), \(c\ge 4\) & \(q f_{2} h_c\). \\
\addlinespace[2pt]
Odd prefix (1, 2, 3, 4, t2), at least five arms & \(q f_{2} h_{t_2} h_{3} f_{4} times later pairs\). \\
\addlinespace[2pt]
Other odd prefix (1, 2, c), at least five arms & \(q f_{2} h_c times later pairs\). \\
\end{longtable}
\endgroup


In the odd \(minimum\ge 2\) row, a denotes the terminal minimum; the
singleton arm length is a separate parameter. The first pair in a five-or-more-arm
unit prefix is displayed explicitly; later pairs are the consecutive
strict residual pairs left after those displayed arms. The rows are
exhaustive by parity, minimum, next integer and the indicated length
exception. The exact unit123 exception is L1 for this reference role;
it is not C106's different ordinary-eighth L2 exception.

Fully reverse this reference construction, including the actual arm,
the L connector, every terminal arm, and both equal-pair decorations.
Superscript minus means precisely this full reversal. Put
\fitdisplay{
X_0=A^L(E_t^-F_{\rm ref}^-),\qquad X=E_r^-X_0,
}{}
\fitdisplay{
Q_a(d)=\begin{cases}1,&a=2,\\1-d/3000,&a\ge3.\end{cases}
}{}
The reference theorem and \(0<E_r^-\le 1\) give
\fitdisplay{
E[\bar f_aX^4]\le E[E_r^-\bar f_aX_0^4]\le Q_a(d)\le1.
}{R45.5.1}
In particular, \(qf_b\) becomes \(p\bar f_b\), and the short \(f_3\bar f_b\) row becomes
\(\bar f_3f_b\). The actual singleton arm points toward its internal root and
has message \(\bar f_a.\) The inner connector points from that root toward its
terminal. The outer connector has \(M=3\) edges from the root toward the
central receiver. Its message is exactly
\fitdisplay{
U=T^3(\bar f_aX).
}{R45.5.2}
The description retains every original edge in its displayed operator order.

\paragraph{R45.5.1 The weighted bridge and positive interpolation}

For every real X and every M, actual row Cauchy gives
\fitdisplay{
|T^M(\bar f_aX)|^2\le \bar f_{M+a}\,T^M(\bar f_a|X|^2).
}{}
Original-measure Holder, then ordinary L2 norm factorization, yield
\fitdisplay{
\begin{aligned}
 I_q(T^M(\bar f_aX))^3
 &\le B_{M+a}\{E[T^M(\bar f_a|X|^2)^2]\}^2\\
 &\le B_{M+a}\|T^M M_{\sqrt{\bar f_a}}\|^4
                  (E[\bar f_a|X|^4])^2.
\end{aligned}
}{R45.5.3}
The vector inside the norm is
\((T^M M_{\sqrt{\bar f_a}})(\sqrt{\bar f_a}|X|^2)\). This is the weighted input and its
matching operator, not an ordinary-fourth substitution into a different
multiplier theorem.

For any positive original operator P and \(w\ge 0,\allowbreak \) row Cauchy gives
\(|P(\sqrt{w}v)|^2\le (P|v|)(P(w|v|)).\) Integration and L2 Cauchy therefore
prove
\fitdisplay{
\|P M_{\sqrt w}\|^2\le\|P\|\,\|P M_w\|.
}{R45.5.4}
Use \(P=T^3,\allowbreak  w=\bar f_a.\) The full-multiplier theorem is
||T \(M_{\bar f_a}||^2\le 1-d/72\) for every \(a\ge 2.\) Factor it in the displayed order:
\(T^3 M_{\bar f_a}=T^2(T M_{\bar f_a}).\) The separate original norm bounds are
\fitdisplay{
\|T^3\|\le1-d/8,\qquad
 \|T^2\|\le(1+(1024/855)\delta_2)^2.
}{}
Squaring (R45.5.4) consequently gives
\fitdisplay{
\|T^3 M_{\sqrt{\bar f_a}}\|^4
 \le(1-d/8)^2(1+(1024/855)\delta_2)^4(1-d/72).
}{R45.5.5}
Set \(c=416/855,\allowbreak \) so (\(1024/855)\delta_2\le c\) d and \(12/25<c<1/2.\)
For the actual \(n=a+3\ge 5,\allowbreak \) (R45.5.1)-(R45.5.5) imply
\fitdisplay{
I_q(U)^3\le
 B_n(1-d/8)^2(1+c d)^4(1-d/72)Q_a(d)^2.
}{R45.5.6}
We keep the genuine \(\delta_2\) until this justified reduction. The remaining
proof pays exactly (R45.5.6), with \(Q_a\le 1.\)

\paragraph{R45.6 Complete continuous payment for every M3 output}

The basic logarithmic inequalities, valid on the indicated positive
domains, are
\fitdisplay{
\log(1-u)\le-u-u^2/2\quad(0\le u<1),
}{}
\fitdisplay{
\log(1+z)\le z-\frac{z^2}{2(1+z)}\quad(z\ge0).
}{R45.6.1}
The second follows by writing \(z-\log(1+z)=\int _0^z\) t/(1+t)dt
and replacing the denominator by 1+z. Since \(d\le 1/3\) and \(c<1/2,\allowbreak \)
\fitdisplay{
4\log(1+c d)\le4c d-\frac38d^2.
}{R45.6.2}
Indeed the exact quadratic coefficient is at least
\(2c^2/(1+c/3)>(12/7)(12/25)^2=1728/4375>3/8.\)

\paragraph{R45.6.1 Domain \(0<d\le 2/15\)}

Use \(B_n\le 1-7d/4\) and the linear logarithmic bounds. The logarithm of
the right side of (R45.5.6), with \(Q_a\le 1,\allowbreak \) is at most
\fitdisplay{
(-7/4-1/4-1/72+4c)d\le-d/72\le-d/80.
}{R45.6.3}
This uses \(4c<2.\) All factors inside the logarithms are positive.

\paragraph{R45.6.2 Domain \(2/15\le d\le 7/40\)}

Use \(B_n\le 1-3d/2,\allowbreak \) retain its quadratic term \(9d^2/8,\allowbreak \) and use (R45.6.2).
The linear coefficient satisfies
\fitdisplay{
4c-7/4-1/72<3/16,
}{}
since \(1664*144=239616<240255=281*855.\) The logarithm is at most
\fitdisplay{
\frac3{16}d-\frac32d^2
 \le(3/16-1/5)d=-d/80.
}{R45.6.4}

\paragraph{R45.6.3 Domain \(7/40\le d\le 1/5\)}

Use \(B_n\le 1-29d/20.\) Its quadratic contribution is \(841d^2/800.\)
Together with (R45.6.2) this gives total quadratic loss \(1141d^2/800.\)
The linear coefficient satisfies
\fitdisplay{
4c-1/4-1/72-29/20<15/64,
}{}
because \(1/4+1/72+29/20=617/360\) and
\(1664*2880=4792320<4797405=5611*855.\) Thus the logarithm is at most
\fitdisplay{
\left(\frac{15}{64}-\frac{1141}{800}\frac7{40}\right)d
 =-\frac{487}{32000}d\le-d/80.
}{R45.6.5}

\paragraph{R45.6.4 Domain \(1/5\le d\le 1/3\)}

Here use (R45.1.3) and \(c<1/2\) in (R45.5.6), and discard only factors already
bounded above by one. The right side is at most
\fitdisplay{
F_*(d)=8(1-2d)(1-5d/2)^3(1+d/2)^4.
}{R45.6.6}
Every factor is positive. Its logarithmic derivative is
\fitdisplay{
-\frac2{1-2d}-\frac{15/2}{1-5d/2}+\frac2{1+d/2}<0.
}{}
The first and last terms have nonpositive sum, and the middle term is
negative. Hence
\fitdisplay{
F_*(d)\le F_*(1/5)=\frac{43923}{50000}<\frac9{10}
 <\frac{239}{240}\le1-d/80\le e^{-d/80}.
}{R45.6.7}
The four domains are exhaustive. They all use the same X, the same
outer \(T^3,\allowbreak \) the same singleton arm, and the same original measure.
This proves Theorem B for every \(d>0.\)

At \(d=0,\allowbreak  K1=0\) on positive support; it does not follow that \(K=0.\)
Original leaf integration gives every actual tree message one, so
\(U=q=1\) and \(Iq(U)^3=1.\) Thus the zero endpoint is also proved.
